\begin{table}[t]
  \centering
  \caption{Useful throughput per format, Nangate45 \SI{45}{\nano\meter}.
           One FLOP is one multiply or one add.}
  \label{tab:residue}
  \setlength{\tabcolsep}{2.2pt}
  \renewcommand{\arraystretch}{1.15}
  \resizebox{\columnwidth}{!}{%
  \begin{tabular}{lrrrrr}
    \toprule
    Format & Elem./ & Instr./ & $F_{\max}$ & Throughput & vs.\ \\
           & instr. & pair    & (\si{\mega\hertz}) & (\si{\giga\flop\per\second}) & \fpeight \\
    \midrule
    \fpfour   & 16 & 1 & 414.9 & 13.28 & 2.88$\times$ \\
    \mtwo     & 16 & 1 & 359.5 & 11.50 & 2.49$\times$ \\
    \fpeight  &  8 & 1 & 288.3 &  4.61 & 1.00$\times$ \\
    \midrule
    \fpfourres & 16 & 2 & 441.0 & 14.11 & 3.06$\times$ \\
    \bottomrule
  \end{tabular}%
  }
  \\[2pt]
  \begin{minipage}{\columnwidth}
    \footnotesize
    Both \insn{mxdotp} and \insn{mxfinal} are structurally pipelined (\S\ref{sec:arch-isa}) with no combinational hazards, a single global stall, and decoupled by an asynchronous mailbox, matching the discipline used by the other three fused engines. The pair's shared clock is bounded by \insn{mxfinal} (\SI{441.0}{\mega\hertz} vs.\ \insn{mxdotp}'s \SI{915.5}{\mega\hertz}). Throughput assumes 2 cycles per completed pair, following directly from this structure.
  \end{minipage}
\end{table}
